\chapter{Analiza și Proiectarea Sistemului Sentinel}
\label{chap:ch3}

\par Proiectarea sistemului Sentinel pornește de la necesitatea de a detecta și neutraliza amenințările cibernetice în timp real. Arhitectura propusă este modulară, decuplând logic instrumentarea nucleului (kernel), evaluarea scorului de risc și aplicarea măsurilor de izolare (containment). Cele trei componente majore interacționează printr-un flux de date unidirecțional: programul eBPF generează evenimente brute, daemonul Go le procesează și le persistă, iar managerul de carantinare le consumă pentru a aplica măsuri de restricție.


\section{Cerințe funcționale și nefuncționale}
\label{sec:requirements}

\par Cerințele funcționale principale care au ghidat proiectarea sunt:

\begin{itemize}
    \item \textbf{RF1}: Interceptarea la nivel de kernel a apelurilor de sistem relevante din perspectiva securității, fără a necesita recompilarea nucleului sau instalarea unor module terțe.
    \item \textbf{RF2}: Calcularea unui scor de risc compus pentru fiecare eveniment detectat, fundamentat pe cel puțin un set de reguli euristice statice.
    \item \textbf{RF3}: Clasificarea decizională a evenimentelor în șase stări de severitate, variind de la activitate benignă (OK) până la necesitatea blocării imediate a procesului (FREEZE).
    \item \textbf{RF4}: Aplicarea automată a acțiunilor de carantinare, strict proporționale cu severitatea detectată, utilizând mecanismele native și standardizate ale kernel-ului Linux.
    \item \textbf{RF5}: Persistarea evenimentelor și a acțiunilor de carantinare într-o bază de date locală, optimizată pentru audit și analiză retrospectivă.
    \item \textbf{RF6}: Expunerea unui serviciu REST API pentru integrarea opțională a unui modul de Machine Learning (ML), capabil de reantrenare dinamică fără întreruperea monitorizării.
\end{itemize}

\par Cerințele nefuncționale critice pentru un mediu de producție sunt:

\begin{itemize}
    \item \textbf{RN1}: Impactul operațional (overhead-ul) introdus de instrumentarea eBPF asupra procesorului (CPU) să fie limitat sub \textbf{5\%} într-un scenariu de activitate normală a sistemului.
    \item \textbf{RN2}: Latența de reacție --- măsurată de la momentul intercepției până la aplicarea efectivă a măsurii de carantinare (excluzând evaluarea ML) --- să nu depășească \textbf{100 ms}.
    \item \textbf{RN3}: Monitorizarea și raportarea continuă a eventualelor pierderi de evenimente cauzate de saturarea memoriei partajate (ring buffer).
\end{itemize}


\section{Arhitectura sistemului}
\label{sec:architecture}

\par Arhitectura sistemului Sentinel se bazează pe trei componente interdependente, ilustrate în schema fluxului de date de mai jos:

\begin{enumerate}
    \item \textbf{Programul eBPF} (\texttt{tracer.bpf.o}): rulează în spațiul privilegiat al nucleului, interceptează apelurile de sistem folosind sonde dinamice (\textit{kprobes}) și scrie evenimentele contextuale într-un \textit{ring buffer} eBPF de înaltă performanță.
    \item \textbf{Daemonul Go} (\texttt{sentinel-daemon}): operează în spațiul utilizator (user-space), consumă datele din ring buffer, evaluează riscul (prin motorul de scoring) și scrie alertele rezultate într-o bază de date SQLite (configurată în modul WAL).
    \item \textbf{Managerul de carantinare} (\texttt{sentinel-quarantine}): interoghează continuu baza de date pentru evenimente cu severitate critică și aplică direct acțiunile de izolare folosind \textit{cgroup v2} și \textit{nftables}.
\end{enumerate}

\begin{figure}[htbp]
\centering
\begin{verbatim}
  Kernel space                          User space
+------------------+   ring buffer   +-------------------+
|   eBPF kprobes   | --------------> |  sentinel-daemon  |
|  (tracer.bpf.o)  |                 | (scoring engine)  |
+------------------+                 +--------+----------+
                                              |  SQLite WAL
                                     +--------v----------+
[optional] <--- ML Scorer REST :8765 |sentinel-quarantine|
                                     | cgroup v2/nftables|
                                     +-------------------+
\end{verbatim}
\caption{Arhitectura sistemului Sentinel: fluxul de date de la kprobe la carantinare.}
\label{fig:architecture}
\end{figure}

\begin{comment}
TODO: check rellevance
\par Decuplarea daemonului de managerul de carantinare, realizată prin intermediul bazei de date SQLite, oferă un avantaj major de reziliență: permite reluarea procesării evenimentelor din punctul ultimei opriri în cazul unui crash (toleranță la defecte) și elimină dependența de sincronizare strictă între producerea evenimentului și aplicarea pedepsei. O latență de interogare asincronă de 500 ms a bazei de date este considerată un compromis acceptabil pentru majoritatea acțiunilor de izolare, menținând în același timp un impact minim asupra I/O-ului sistemului.
\end{comment}

\section{Proiectarea programului eBPF}
\label{sec:ebpf_design}

\par Componenta eBPF este proiectată să asigure interceptarea a 20 de apeluri de sistem relevante. Fiecare apel este instrumentat printr-un punct de interceptare kprobe dedicată, având prefixul \texttt{kp\_}. Arhitectura folosește două tipologii principale de rutine de tratare (handlers):

\begin{itemize}
    \item \texttt{submit\_event(nr)}: funcție generică utilizată pentru apelurile care nu operează cu căi de fișiere. Extrage un set esențial de metadate: marca temporală (nanosecunde), PID, UID, numărul apelului și denumirea executabilului (\texttt{comm}).
    \item \texttt{submit\_event\_with\_path(nr, ctx)}: variantă extinsă, utilizată pentru apelurile care manipulează fișiere (ex. \texttt{openat}, \texttt{unlinkat}, \texttt{renameat2}). Această rutină utilizează asistentul \texttt{bpf\_probe\_read\_user\_str} pentru a extrage dinamic pointer-ul către șirul de caractere (din registrul \texttt{rsi} al structurii \texttt{pt\_regs}) peste granița spațiului utilizator.
\end{itemize}

\par Handler-ul pentru \texttt{write} (syscall 1) necesită un tratament special. Deoarece argumentul său este un descriptor de fișier (număr întreg), nu o cale, handler-ul eBPF ignoră explicit scrierile către descriptorii standard (0, 1 și 2) --- acestea generează un volum ridicat de evenimente cu relevanță de securitate scăzută. Pentru descriptorii rămași, funcția codifică valoarea numerică într-un șir de forma \texttt{"fd:N"}. Sarcina traducerii este delegată daemonului Go din spațiul utilizator, care interoghează pseudo-sistemul \texttt{/proc/\{pid\}/fd/\{N\}} pentru a obține calea absolută, menținând astfel execuția eBPF fluidă \cite{Kerrisk2010}.

\par Filtrul pe tabela de dispersie \texttt{pid\_blacklist} previne fenomenul de auto-instrumentare (detectia propriilor apeluri de sistem): imediat după lansare, daemonul Go își scrie propriul PID în această structură partajată de tip dicționar. La fiecare kprobe, codul eBPF verifică această listă și abandonează silențios interceptările generate de componentele de securitate proprii, eliminând astfel bucla de feedback ce ar satura ring buffer-ul.


\section{Proiectarea motorului de scoring}
\label{sec:scoring_design}

\par Motorul de scoring rafinează progresiv nivelul de amenințare printr-un sistem decizional etapizat, generând un scor final în intervalul $[0, 100]$.

\subsection{Etapa 1: Evaluarea statică a apelurilor de sistem}
\label{subsec:static_weights}

\par Fiecărui apel interceptat îi este asociată o pondere euristică $w_{nr} \in [5, 95]$, calibrată empiric conform studiilor privind vectorii de atac Linux \cite{Cozzi2018, Buczak2016}. Ponderile extreme (95) sunt alocate acțiunilor care nu au utilitate legitimă normală (ex. \texttt{init\_module}), în timp ce valorile minime (5) corespund apelurilor zgomotoase (\texttt{openat}). Dacă procesul inițiator beneficiază de privilegii administrative ($uid = 0$), ponderea este majorată cu un factor de 1,3, reflectând gravitatea exponențială a atacurilor executate cu drepturi depline:

\begin{equation}
    \text{score}_{\text{static}}(nr, uid) = w_{nr} \cdot \begin{cases} 1.3 & \text{dacă } uid = 0 \\ 1.0 & \text{altfel} \end{cases}
    \label{eq:static_score}
\end{equation}

\par Amplificarea pentru procesele root reflectă faptul că aceleași syscall-uri au un impact semnificativ mai mare când sunt executate cu privilegii complete.


\subsection{Etapa 2: Analiza comportamentală a activității în rafală (Burst Detection)}
\label{subsec:burst_detection}

\par Atacurile de tip forță brută (SSH brute-force) sau epuizare a resurselor (fork bombs) se manifestă prin frecvențe anormale de apeluri identice. Etapa 2 implementează o fereastră temporală glisantă (sliding window) care monitorizează agregarea pe tripletul (\texttt{PID}, \texttt{syscall\_nr}, \texttt{start\_time}). Includerea timpului de inițializare a procesului (\texttt{start\_time}) previne coliziunile cauzate de fenomenul de \textit{PID wrap-around} în sistemele cu rotație rapidă a proceselor.

\begin{table}[htbp]
\begin{center}
\begin{tabular}{|p{90pt}|p{50pt}|p{70pt}|p{55pt}|p{60pt}|}
\hline
\textbf{Syscall(uri)} & \textbf{Prag} & \textbf{Fereastră} & \textbf{Cooldown} & \textbf{Scor burst} \\
\hline\hline
\texttt{socket}, \texttt{connect} & 5  & 3 s  & 10 s & 75.0 \\
\hline
\texttt{clone}                     & 20 & 2 s  & 10 s & 35.0 \\
\hline
\texttt{openat}                    & 8  & 3 s  & 10 s & 60.0 \\
\hline
\end{tabular}
\end{center}
\caption{Parametrii de detectare a activității în rafale}
\label{tab:burst_params}
\end{table}

\par Intervalul de cooldown de 10 s previne fenomenul de \textit{alert fatigue} --- odată emis un alert pentru un triplet (PID, syscall), acesta nu va mai fi reeditat timp de 10 secunde pentru aceeași manifestare continuă, reducând zgomotul în baza de date.


\subsection{Etapa 3: Analiza contextuală bazată pe căile fișierelor}
\label{subsec:path_alerts}

\par A treia etapaă realizează o potrivire de tipare (pattern matching) pe căile fișierelor accesate de syscall-urile relevante (\texttt{write} cu fd rezolvat, \texttt{openat}, \texttt{unlinkat}, \texttt{renameat2}), vizând tactici cunoscute de post-exploatare:

\begin{itemize}
    \item \textbf{Log Tamper} (scor 50): scriere sau alterare a directoarelor \texttt{/var/log/}, specifică tentativelor de ștergere a urmelor.
    \item \textbf{LSMFileWrite} (scor 50): manipularea fișierelor \texttt{/etc/cron*} (asigurarea persistenței) sau deschiderea bibliotecilor dinamice (\texttt{.so}) din directoare mapate în memorie (\texttt{/tmp/}, \texttt{/dev/shm/}), un indicator pentru executarea \textit{fileless malware}.
    \item \textbf{Privilege Escalation} (scor 65): alterarea directă a bazelor de date de autentificare (\texttt{/etc/passwd}, \texttt{/etc/shadow}, \texttt{/etc/sudoers}).
\end{itemize}

\par Dacă o regulă contextuală este validată, scorul acesteia \textbf{suprascrie} scorul static inițial (în loc să se cumuleze prin adunare), prevenind astfel inflația artificială a riscului și reflectând cu acuratețe severitatea tacticii interceptate.


\subsection{Integrarea scorerului ML}
\label{subsec:ml_integration}

\par Dacă URL-ul scorerului ML este configurat la pornire, daemonul Go apelează asincron endpoint-ul \texttt{POST /score} pentru fiecare eveniment cu scor $\geq 10$. Apelul este condiționat de o limită (\textit{timeout}) de 50 ms; în absența unui răspuns în timp util, evaluarea continuă exclusiv static. Scorul final ia forma unei fuziuni ponderate (Score Blending):

\begin{equation}
    \text{score}_{\text{final}} = (1 - w) \cdot \text{score}_{\text{static}} + w \cdot \text{score}_{\text{ML}}
    \label{eq:blended_score}
\end{equation}

\noindent unde $w \in \{0.1, 0.2, 0.5, 0.7\}$ variază invers proporțional cu gravitatea scorului static:

\begin{itemize}
    \item $w_{nr} \geq 85$: $w = 0.1$ --- pentru vectori destructivi clari, logica statică domină absolut decizia (90\%).
    \item $w_{nr} \in [60, 85)$: $w = 0.2$ --- analiza statică păstrează o autoritate de 80\%.
    \item $w_{nr} \in [30, 60)$: $w = 0.5$ --- decizie echilibrată (50/50) între reguli și comportamentul algoritmic.
    \item $w_{nr} < 30$: $w = 0.7$ --- sistemul ML primește libertate majoritară de a discerne anomaliile fine în zgomotul apelurilor comune.
\end{itemize}

\par Un floor de $0.5 \cdot \text{score}_{\text{static}}$ garantează că modulul ML nu poate submina evaluarea de securitate la mai puțin de jumătate din valoarea sa originară, prevenind atacatorii să camufleze exploit-uri severe în tipare artificiale de \textit{trafic normal}.


\section{Mașina de stări pentru carantinare}
\label{sec:state_machine}

\par Severitatea finală este tradusă operațional într-una din cele șase stări acționabile, ordonate crescător: OK $\rightarrow$ WATCH $\rightarrow$ WARN $\rightarrow$ THROTTLE $\rightarrow$ ISOLATE $\rightarrow$ FREEZE, conform pragurilor din Tabelul~\ref{tab:states}.

\begin{table}[htbp]
\begin{center}
\begin{tabular}{|p{65pt}|p{70pt}|p{175pt}|}
\hline
\textbf{Stare} & \textbf{Prag scor} & \textbf{Acțiune aplicată} \\
\hline\hline
OK       & $< 30$    & Eveniment considerat benign și filtrat; nicio acțiune \\
\hline
WATCH    & $\geq 30$ & Eveniment suspect, jurnalizat pentru auditare în baza de date \\
\hline
WARN     & $\geq 50$ & Alertare activă; eveniment propagat în dashboard \\
\hline
THROTTLE & $\geq 70$ & Limitare CPU la 10\% prin directiva \texttt{cpu.max} în cgroup v2 \\
\hline
ISOLATE  & $\geq 85$ & Limitare de memorie la 256 MiB și izolare completă a rețelei via nftables \\
\hline
FREEZE   & $\geq 95$ & Suspendarea procesului la nivel de kernel (\texttt{SIGSTOP} și \texttt{cgroup.freeze}); autorecuperare programată după 5 minute \\
\hline
\end{tabular}
\end{center}
\caption{Mașina de stări: praguri de scor și acțiuni de carantinare corespunzătoare}
\label{tab:states}
\end{table}

\par Aceste praguri au fost aliniate operațional: un apel banal executat de contul de sistem (\texttt{openat} ca root, scor $6.5$) este ignorat; o rafală nejustificată de conexiuni atinge nivelul THROTTLE; în timp ce simpla încercare de instalare a unui modul extern (\texttt{finit\_module} ca root, scor $123.5$, plafonat la $95$) atrage suspendarea absolută (FREEZE) a procesului atacator.

\par Izolarea rețelei în starea ISOLATE se realizează prin nftables: managerul de carantinare creează un tabel dedicat \texttt{inet sentinel\_quarantine} și aplică reguli restrictive de filtrare (drop) bazate pe identificatorul (inode-ul) atribuit de cgroup procesului compromis \cite{Menage2007}. Această granulare asigură restricționarea chirurgicală a ramurii malițioase, fără a afecta comunicațiile serverului de producție.


\section{Schema bazei de date}
\label{sec:db_schema}

\par Sentinelul utilizează SQLite în modul WAL (Write-Ahead Logging) cu parametrul activ de gestionare a blocajelor (\texttt{\_busy\_timeout=5000}). Această configurație tranzacțională suportă citiri recurente (efectuate de managerul de carantinare) care nu invalidează și nu blochează fluxul continuu de scriere asigurat de daemonul principal. SQLite WAL permite citiri concurente cu o scriere activă, satisfăcând astfel modelul de acces al sistemului.

\par Sunt menținute două tabele principale:

\begin{itemize}
    \item \texttt{events}: jurnalizează totalitatea evenimentelor relevante. Câmpurile includ \texttt{alert\_type}, \texttt{score}, \texttt{state}, \texttt{pid}, \texttt{comm}, \texttt{timestamp} (ISO 8601) și un payload flexibil de metadate \texttt{detail} (în format JSON, conținând câmpurile \texttt{ts\_ns}, \texttt{syscall\_name}, \texttt{fname} și \texttt{source}).
    \item \texttt{quarantine\_log}: o pistă de audit irefutabilă pentru măsurile de remediere automată. Leagă acțiunea de izolare de declanșatorul inițial prin cheia externă \texttt{event\_id} și înregistrează metadatele restricției aplicate: acțiunile sistemice (\texttt{actions} codificat JSON), confirmarea succesului, erorile OS și durata temporală de execuție a pedepsei (\texttt{latency\_ms}).
\end{itemize}
 
\par Câmpul \texttt{latency\_ms} din \texttt{quarantine\_log} înregistrează durata totală de aplicare a acțiunilor de containment, de la selectarea evenimentului până la confirmarea acțiunilor, oferind un mecanism integrat prin care performanța de reacție a sistemului poate fi evaluată și validată obiectiv în concordanță cu cerința nefuncțională RN2.